\begin{center}
    {\LARGE \textbf{2009无机化学}} \\[2ex]
    {\large 时量180分钟 \quad 满分100分}
\end{center}

\vspace{0.5cm}
\noindent\textbf{注意事项:}
\begin{enumerate}[label=\arabic*., parsep=0pt, itemsep=0pt]
    \item 答题前,考生先将自己的姓名、学号、考场号、座位号填写在答题卡上,将条形码准确粘贴在考生信息条形码粘贴区。
    \item 考试结束后,将本试卷与答题纸一并收回。
    \item 回答各个问题时,将答案写在答题纸上,写在本试卷上无效。
\end{enumerate}

\vspace{0.5cm}
\noindent\textbf{一、完成并配平下列反应的方程式（20 pts.）}

\begin{enumerate}[label=\arabic*., start=1, itemsep=1.5ex]
    \item 用氢碘酸溶液处理\ce{Fe2O3}。
    \item 向\ce{KMnO4}溶液中通入过量的\ce{SO2}。
    \item 用酸性\ce{H2O2}处理\ce{MnO2}。
    \item 向溴水中通入\ce{H2S}。
    \item 向\ce{HgCl2}溶液中滴加\ce{SnCl2}溶液至过量。
    \item 向\ce{CrCl3}溶液中加入\ce{NaOH}溶液至过量。
\end{enumerate}

\vspace{0.5cm}
\noindent\textbf{二、解释下列实验现象（28 pts.）}

\begin{enumerate}[label=\arabic*., start=7, itemsep=1.5ex]
    \item 向\ce{Hg(NO3)2}溶液中加入\ce{NaOH}溶液生成黄色沉淀，向\ce{Hg(NO3)2}溶液中加入氨水则生成白色沉淀。
    \item 将\ce{SnCl2}溶液倒入大量水中对物质：
    \item 向\ce{CuSO4}溶液中滴加\ce{KI}溶液，有黄色沉淀生成，再滴加过量\ce{Na2S2O3}溶液，则沉淀先变白而后消失，得到无色溶液。
    \item 向\ce{FeSCN^{2+}}溶液中滴加盐酸，红色不消失；向\ce{FeSCN^{2+}}溶液中滴加氢氟酸，则红色消失。
\end{enumerate}

\vspace{0.5cm}
\noindent\textbf{三、简要回答下列各题（52 pts.）}

\begin{enumerate}[label=\arabic*., start=11, itemsep=1.5ex]
    \item 给出原子序数为24、44元素的名称、符号和价层电子排布式。
    \item 给出下列分子或离子的几何构型和中心原子的杂化类型
    \begin{enumerate}[label=(\arabic*), itemsep=1ex]
        \item \ce{NO2}
        \item \ce{SOCl2}
        \item \ce{XeF4}
        \item \ce{[PtCl4]^{2-}}
    \end{enumerate}
    \item 将\ce{KMnO4}加入用硫酸酸化的\ce{KBr}溶液时，\ce{KMnO4}立即退色；将\ce{KMnO4}加入用醋酸酸化的\ce{KBr}溶液时，\ce{KMnO4}退色很慢。请解释。
    \item 请解释：\ce{CuSO4}为白色，\ce{CuCl2 * 2H2O}为绿色，\ce{CuSO4 * 5H2O}为蓝色。
    \item 请解释：\ce{NO}和\ce{B2}为顺磁性物质，\ce{NO^{+}}和\ce{CO}为抗磁性物质。
    \item 比较各对化合物热稳定性并说明原因：
    \begin{enumerate}[label=(\arabic*), itemsep=1ex]
        \item \ce{CuO}与\ce{HgO}
        \item \ce{K2SO3}与\ce{K2SO4}
        \item \ce{K2CO3}与\ce{Li2CO3}
    \end{enumerate}
    \item 比较下列各对配合物的稳定性并说明原因：
    \begin{enumerate}[label=(\arabic*), itemsep=1ex]
        \item \ce{[Cu(NH3)4]^{2+}}与\ce{[Zn(NH3)4]^{2+}}
        \item \ce{[Cu(NH3)4]^{2+}}与\ce{[Cu(NH3)4]^{+}}
        \item \ce{[Ni(CN)4]^{2-}}与\ce{[Pt(CN)4]^{2-}}
        \item \ce{[Zn(SCN)4]^{2-}}与\ce{[Hg(SCN)4]^{2-}}
    \end{enumerate}
    \item 解释：常温下\ce{MnO2}为固体而\ce{Mn2O7}为液体。
\end{enumerate}

\vspace{0.5cm}
\noindent\textbf{四、回答下列各题（50 pts.）}

\begin{enumerate}[label=\arabic*., start=19, itemsep=1.5ex]
    \item 用反应方程式表示下列转化过程：\ce{Fe -> FeSO4 -> (NH4)2SO4\ * FeSO4 * 6H2O -> K3[Fe(CN)6]}
    \item 设计方案分离溶液中的\ce{Ag^{+}}、\ce{Ni^{2+}}、\ce{Zn^{2+}}、\ce{Na^{+}}。
    \item 用二种方法鉴别下列各对物质：
    \begin{enumerate}[label=(\arabic*), itemsep=1ex]
        \item \ce{Na2SO3}和\ce{Na2CO3}
        \item \ce{Na2HPO3}和\ce{Na2HPO4}
    \end{enumerate}
    \item 黄色晶体\ce{A}溶于稀盐酸后通入\ce{SO2}得绿色溶液\ce{B}，再加入过量\ce{KOH}溶液得绿色溶液\ce{C}。溶液\ce{C}中加\ce{H2O2}得黄色溶液\ce{D}。\ce{A}的溶液中加入硝酸铅溶液生成黄色沉淀\ce{E}。\ce{A}的溶液中加入硫酸和\ce{H2O2}有蓝色的\ce{F}生成。溶液\ce{B}中入\ce{Na2CO3}有沉淀\ce{G}生成。干燥的\ce{G}与\ce{KClO3}和\ce{KOH}共熔又有\ce{A}生成。请给出\ce{A} $\sim$\ ce{G}的化学式。
    \item 计算由如下电极组成原电池的电动势。
    \begin{itemize}
        \item 电极A：铜片插入$2\ \mathrm{mol\cdot dm^{-3}}$\ \ce{CuSO4}和$10\ \mathrm{mol\cdot dm^{-3}}$氨水的混合溶液；
        \item 电极B：银-氯化银标准电极；
    \end{itemize}
    已知：
    \[
    \begin{array}{ll}
    \ce{Cu^{2+} + 2e^{-} = Cu} & E^\ominus = 0.34\ \mathrm{V} \\
    \ce{Ag^{+} + e^{-} = Ag} & E^\ominus = 0.80\ \mathrm{V} \\
    \ce{Cu^{2+} + 4NH3 = [Cu(NH3)4]^{2+}} & K_{\text{稳}} = 10^{12} \\
    \ce{Ag^{+} + Cl^{-} = AgCl} & K_{\text{sp}} = 10^{-10} \\
    \end{array}
    \]
    \item 试解释：键角\ce{ClO2}（$118^\circ$）大于\ce{Cl2O}（$111^\circ$），而键长\ce{ClO2}小于\ce{Cl2O}。
\end{enumerate}